% Generated by scripts/assemble.py; corrections belong in source/corrections.json.
\examhead{134}{华中科技大学 2026 年数学分析试卷(回忆版)}{128}
% p129-b01; source PDF page 129
\blockanchor{p129-b01}
(注.题目来自考生回忆,部分表述或数据可能有误)\par
% p129-b02; source PDF page 129
\blockanchor{p129-b02}
\question{1}(10 分) 求极限 $\lim_{n\to \infty}\sum_{k = 1}^{n}\sin \frac{2k - 1}{n^2}.$\par
% p129-b03; source PDF page 129
\blockanchor{p129-b03}
\question{2}(10 分) 求极限 $\lim_{x\to 0}\frac{\tan(\tan x) - \sin(\sin x)}{(2026^x - 1)\ln(1 + x^2)}.$\par
% p129-b04; source PDF page 129
\blockanchor{p129-b04}
\question{3}(10 分) 设 $\theta \in (-1, 1)$, 试求 $\int_{0}^{\frac{\pi}{2}} \ln \left( \frac{1 + \theta \cos x}{1 - \theta \cos x} \right) \frac{1}{\cos x} \mathrm{d}x$.\par
% p129-b05; source PDF page 129
\par\Needspace{6\baselineskip}
\blockanchor{p129-b05}
\question{4}(10 分) 设 $f(x) = \frac{1}{\sqrt{4\pi^2 - x^2}}$ 在 $[- \pi, \pi]$ 上的傅里叶级数展开式为\par
\nopagebreak[4]
% p129-b06; source PDF page 129
\blockanchor{p129-b06}
\displayformula{f (x) \sim \frac {a _ {0}}{2} + \sum_ {n = 1} ^ {\infty} \left(a _ {n} \cos n x + b _ {n} \sin n x\right).}
% p129-b07; source PDF page 129
\blockanchor{p129-b07}
求 $\sum_{n=1}^{\infty}(a_n^2 + b_n^2)$.\par
% p129-b08; source PDF page 129
\par\Needspace{6\baselineskip}
\blockanchor{p129-b08}
\question{5}(10 分) 设点 $A(0,0), B(\pi, \pi^{\frac{2}{3}})$, $\Gamma$ 是由 $A$ 沿 $y = x^{\frac{2}{3}}$ 到 $B$ 的曲线段, 试求\par
\nopagebreak[4]
% p129-b09; source PDF page 129
\blockanchor{p129-b09}
\displayformula{\int_ {\Gamma} \frac {\cos y}{\sqrt {x ^ {2} + 1}} \mathrm{d} x - \ln (x + \sqrt {x ^ {2} + 1}) \sin y \mathrm{d} y.}
% p129-b10; source PDF page 129
\blockanchor{p129-b10}
\question{6}(10 分) 区域 $\Omega = \{(x, y, z) \mid x^2 + y^2 \leq 1, 0 \leq z \leq \sqrt{x^2 + y^2}\}$, $S = \partial \Omega$, 正方向向外, 求曲面积分 $\iint_S\cdots$\par
% p129-b11; source PDF page 129
\blockanchor{p129-b11}
(具体数据忘了, 是经典的 Gauss 公式应用问题)\par
% p129-b12; source PDF page 129
\blockanchor{p129-b12}
\question{7}(15 分) 设  $x_{1}=\frac{2025}{2026}$ ,  $x_{n+1}=\sin x_{n}(n=1,2,\cdots)$ .\par
% p129-b13; source PDF page 129
\blockanchor{p129-b13}
(1) 证明: $\lim_{n\to \infty}x_n = 0.$\par
% p129-b14; source PDF page 129
\blockanchor{p129-b14}
(2) 试证明: 存在正整数 $m$, 使得存在 $c \in \mathbb{R}$, 有 $\lim_{n \to \infty} n^{\frac{1}{m}} x_n = c$.\par
% p129-b15; source PDF page 129
\blockanchor{p129-b15}
\question{8}(15 分) $\{a_{n}\}, \{b_{n}\}$ 满足 $\sum_{n=1}^{\infty}|a_{n+1} - a_{n}|$ 收敛, $\lim_{n\to \infty}a_n = 0, \{b_n\}$ 的部分和有界, 证明: $\sum_{n=1}^{\infty}a_{n}b_{n}$ 收敛.\par
% p129-b16; source PDF page 129
\par\Needspace{6\baselineskip}
\blockanchor{p129-b16}
\question{9}(15 分) 设 $f(x)$ 一阶连续可导, 且 $f(0) = f'(0) = 0$, $f''(0)$ 存在, 证明:\par
\nopagebreak[4]
% p129-b17; source PDF page 129
\blockanchor{p129-b17}
\displayformula{\lim _ {x \rightarrow 0} \frac {f (\tan x) - f (\sin x)}{x ^ {3} \ln (1 + x)} = \frac {1}{2} f ^ {\prime \prime} (0).}
% p129-b18; source PDF page 129
\blockanchor{p129-b18}
\question{10}(15 分) 设 $f(x) \in C[a, b]$, 证明: $\lim_{n \to \infty} \frac{1}{n} \ln \left( \int_{a}^{b} e^{-nf(x)} \mathrm{d}x \right) = -\min_{x \in [a, b]} f(x)$.\par
% p129-b19; source PDF page 129
\blockanchor{p129-b19}
\question{11}(15 分) 证明: $\sum_{n=1}^{\infty} \frac{\sin nx}{n} \arctan x$ 在 $[0,1]$ 上一致收敛.\par
% p129-b20; source PDF page 129
\blockanchor{p129-b20}
\question{12}(15 分) 一个巨麻烦的积分不等式证明.\par
